\begin{figure}
\caption{Resultados do desempenho F1 dos circuitos 34, 35, 37 e 38 com os modelos geradores de \textit{embeddings} na base \textit{ChatGPT} - Fácil. Esses modelos implementam modelos clássicos não-lineares. As barras comparativas do lado direito exibem a diferença percentual da métrica F1 entre o modelo com a maior média F1 e o modelo \textit{Matryoshka}, divididos pela profundidade de cada circuito. As barras de diferença são coloridas com a mesma cor do melhor modelo em questão. As marcações que indicam semelhança estatística são advindas do teste pareado de Wilcoxon com correção de p-valor de Holm.}
\label{fig:model_transformer:easy:g5}
     \centering
    \centering
    \resizebox{!}{0.90\textheight}{%  <-- Ajuste a altura aqui
        \import{capitulos/4_resultados/figuras}{transformers_chatgpt_easy_g5.pgf}
    }

      \par\medskip\ABNTEXfontereduzida\selectfont\textbf{Fonte:} Elaborada pelo autor (2025) \par\medskip
\end{figure}